\documentclass{article}
\usepackage[utf8]{inputenc}
\usepackage{a4wide}

\begin{document}

\section*{ศึกสุดท้าย}

สงครามด่านสุดท้ายสำหรับเหล่านักรบ~heroes~กลุ่มหนึ่งที่กำลังจะต่อสู้กับปีศาจmonstersจำนวนหนึ่ง

กำหนดอาร์เรย์เลขจำนวนเต็มบวก (เริ่มดัชนีแรกที่1) ให้สองชุด

ได้แก่\texttt{heroes}เป็นพลังชีวิตของเหล่านักรบnคน

และ\texttt{monsters}เป็นของปีศาจ\texttt{m}ตัว

กำหนดว่านักรบคนที่iสามารถปราบปีศาจตัวที่jได้ก็ต่อเมื่อ\texttt{monsters{[}j{]}\textless{}=\ heroes{[}i{]}}



นอกจากนี้คุณยังได้อาร์เรย์เลขจำนวนเต็มบวก~coin(เริ่มดัชนีแรกที่1)

ความยาวmแทนจำนวนเหรียญ

นักรบที่ปราบปีศาจตัวที่~k~ได้จะได้รับเหรียญ\texttt{coin{[}k{]}}



ให้ตอบมาเป็นอาร์เรย์ansความยาวnโดยที่\texttt{ans{[}p{]}}เป็นเหรียญสูงสุดที่นักรบคนที่pสามารถสะสมได้ในสงครามครั้งสุดท้ายนี้

ทั้งนี้...



1.พลังชีวิตของนักรบจะ\ul{ไม่}ลดลงหลังจากปราบปีศาจแต่ละตัว



2.นักรบหลายคนสามารถปราบปีศาจตัวหนึ่งๆได้

แต่ปีศาจแต่ละตัวจะถูกปราบโดยนักรบแต่ละคนได้เพียงครั้งเดียว



\hfill\break



{\def\LTcaptype{none} % do not increment counter

\begin{longtable}[]{@{}

  >{\raggedright\arraybackslash}p{(\linewidth - 4\tabcolsep) * \real{0.3333}}

  >{\raggedright\arraybackslash}p{(\linewidth - 4\tabcolsep) * \real{0.3333}}

  >{\raggedright\arraybackslash}p{(\linewidth - 4\tabcolsep) * \real{0.3333}}@{}}

\toprule\noalign{}

\begin{minipage}[b]{\linewidth}\centering

ข้อมูลเข้า

\end{minipage} & \begin{minipage}[b]{\linewidth}\centering

ข้อมูลออก

\end{minipage} & \begin{minipage}[b]{\linewidth}\centering

คำอธิบาย

\end{minipage} \\

\midrule\noalign{}

\endhead

\bottomrule\noalign{}

\endlastfoot

\begin{minipage}[t]{\linewidth}\raggedright

3 5\\

1 4 2\\

1 1 5 2 3\\

2 3 4 5 6\\

\strut

\end{minipage} & \begin{minipage}[t]{\linewidth}\raggedright

5 16 10\\

\strut

\end{minipage} & \begin{minipage}[t]{\linewidth}\raggedright

สำหรับฮีโร่แต่ละคน เราจะแจงดัชนีของเหล่าปีศาจที่เข้าปราบได้:\\

\strut \\

ฮีโร่ 1: {[}1,2{]} เนื่องจากพลังของฮีโร่คนนี้คือ 1 เป็น monsters{[}1{]},

monsters{[}2{]} \textless= 1 ดังนั้นฮีโร่คนนี้จึงสะสมเหรียญได้\\

coins{[}1{]} + coins{[}2{]} = 5 เหรียญ\\

\strut \\

ฮีโร่ 2:\\

{[}1,2,4,5{]} เนื่องจากพลังของฮีโร่คนนี้คือ 4\\

และ monsters{[}1{]}, monsters{[}2{]}, monsters{[}4{]}, monsters{[}5{]}

\textless= 4 ดังนั้นฮีโร่คนนี้จะสะสมเหรียญได้ coins{[}1{]} + coins{[}2{]} +

coins{[}4{]} + coins{[}5{]} = 16 เหรียญ\\

\strut \\

ฮีโร่ 3:\\

{[}1,2,4{]} เนื่องจากพลังของฮีโร่คนนี้คือ 2 และ monsters{[}1{]},\\

monsters{[}2{]}, monsters{[}4{]} \textless= 2 ดังนั้นฮีโร่คนนี้สะสมเหรียญได้\\

coins{[}1{]} + coins{[}2{]} + coins{[}4{]} = 10เหรียญ.\\

\strut \\

ดังนั้นคำตอบจะเป็น~ 5 16 10\\

\strut

\end{minipage} \\

\begin{minipage}[t]{\linewidth}\raggedright

1 4\\

5\\

2 3 1 2\\

10 6 5 2\\

\strut

\end{minipage} & 23 & \begin{minipage}[t]{\linewidth}\raggedright

ฮีโร่คนนี้ปราบปีศาจได้ทั้งหมด เนื่องจาก monsters{[}i{]} \textless= 5

ดังนั้นจึงสะสมเหรียญได้ทั้งหมด coins{[}1{]} + coins{[}2{]} + coins{[}3{]}

+coins{[}4{]} = 23 และตอบ 23\\

\strut

\end{minipage} \\

\end{longtable}

}



\subsection*{Input}
บรรทัดที่1เลขจำนวนเต็มบวก เป็นค่าnและmเว้นด้วยช่องว่าง



บรรทัดที่2ชุดเลขจำนวนเต็มบวกnค่าเป็นพลังชีวิตของเหล่านักรบ เว้นด้วยช่องว่าง



บรรทัดที่3ชุดเลขจำนวนเต็มบวกnค่าเป็นพลังชีวิตของปีศาจ เว้นด้วยช่องว่าง



บรรทัดที่4ชุดเลขจำนวนเต็มบวกmค่าเป็นค่าเหรียญ เว้นด้วยช่องว่าง



\subsection*{Output}
ชุดเลขจำนวนเต็มansความยาวnโดยที่ans{[}p{]}เป็นเหรียญสูงสุดที่นักรบคนที่pสามารถสะสมได้\\



\end{document}
